% GENERATED FILE. Edit canonical lessons, then run npm run generate:books.
% canonical-source-hash: fnv1a64:36c8f6a74dff7103
% canonical-lessons: MR-C365-vataghati-karne, MR-C365-vad-ghalne, MR-C365-retne, MR-C365-mitavne, MR-C365-sangharsa-karne, MR-R365-review-358-365

\chapter{Settling Disputes}
\label{ch:mr-settling-disputes}
\hlchaptermodality{365}

\begin{chapteropening}
\textbf{By the end of this chapter:} \emph{I can talk about arguing, negotiating and settling a quarrel.}

Complete the last lesson of chapter 365: I can talk about arguing, negotiating and settling a quarrel.
\end{chapteropening}

\section[\texorpdfstring{\mr{वाटाघाटी} \mr{करणे} (vāṭāghāṭī karṇe)}{वाटाघाटी करणे (vāṭāghāṭī karṇe)}]{\mr{वाटाघाटी} \mr{करणे} (vāṭāghāṭī karṇe) — to negotiate}

\label{lesson:MR-C365-vataghati-karne}



\begin{quote}
Before the new one: say the Marathi for a spokesperson, then the Marathi for an activist.
\end{quote}

\subsection*{You'll want to know: \mr{वाटाघाटी} \mr{करणे}}
\textbf{\mr{वाटाघाटी} \mr{करणे}} — \emph{vāṭāghāṭī karṇe} — \textquotedblleft{}to negotiate\textquotedblright{}.

\begin{quote}
\mr{दोन्ही} \mr{बाजूंनी} \mr{रात्रभर} \mr{वाटाघाटी} \mr{केल्या}.
\end{quote}

\begin{quote}
\emph{donhī bājūnnī rātrabhar vāṭāghāṭī kelyā.}
\end{quote}

\begin{quote}
Both sides negotiated all night.
\end{quote}

Past, with -\mr{ने} (he …): \mr{त्याने} \mr{वाटाघाटी} \mr{केल्या}, \emph{tyāne vāṭāghāṭī kelyā}.

\mr{वाटाघाटी} is feminine plural, so the past is \mr{केल्या}.

\begin{cousinweb}
\textbf{\mr{वाटाघाटी} \mr{करणे}} (\emph{vāṭāghāṭī karṇe}) joins \mr{वाटाघाटी}, \textquotedblleft{}talks, bargaining\textquotedblright{}, a Marathi word whose origin is not settled, to \mr{करणे}, \textquotedblleft{}to do\textquotedblright{}.
\end{cousinweb}

\begin{grammarlens}[title={What you've built}]
One more word for giving reasons and opinions.
\end{grammarlens}

\subsection*{Your turn}
\begin{itemize}
\raggedright
  \item \emph{Say it:} \emph{vāṭāghāṭī karṇe}
  \item \emph{Say it:} \emph{vāṭāghāṭī karṇe}, once more
  \item \emph{Say it:} \emph{donhī bājūnnī rātrabhar vāṭāghāṭī kelyā.} — Both sides negotiated all night.
  \item \emph{Recall:} say the Marathi for a spokesperson, then the Marathi for an activist, then say \emph{vāṭāghāṭī karṇe} again
\end{itemize}

\begin{tcolorbox}[breakable,colback=teal!4,colframe=teal!35!black,title={Before you move on}]
What does \mr{वाटाघाटी} \mr{करणे} mean? (\textquotedblleft{}To negotiate\textquotedblright{}.) Say it once more.
\end{tcolorbox}
\section[\texorpdfstring{\mr{वाद} \mr{घालणे} (vād ghālṇe)}{वाद घालणे (vād ghālṇe)}]{\mr{वाद} \mr{घालणे} (vād ghālṇe) — to argue}

\label{lesson:MR-C365-vad-ghalne}



\begin{quote}
Before the new one: say the Marathi for an activist, then the Marathi for to negotiate.
\end{quote}

\subsection*{You'll want to know: \mr{वाद} \mr{घालणे}}
\textbf{\mr{वाद} \mr{घालणे}} — \emph{vād ghālṇe} — \textquotedblleft{}to argue\textquotedblright{}.

\begin{quote}
\mr{मुलं} \mr{खेळण्यांवरून} \mr{वाद} \mr{घालत} \mr{होती}.
\end{quote}

\begin{quote}
\emph{mulaṁ kheḷaṇyāṁvarūn vād ghālat hotī.}
\end{quote}

\begin{quote}
The children were arguing over the toys.
\end{quote}

Past, with -\mr{ने} (he …): \mr{त्याने} \mr{वाद} \mr{घातला}, \emph{tyāne vād ghātlā}.

The subject of the quarrel takes -\mr{वरून}; \mr{घालणे}'s past is \mr{घातला}.

\begin{cousinweb}
\textbf{\mr{वाद} \mr{घालणे}} (\emph{vād ghālṇe}) joins the tatsama \mr{वाद}, Sanskrit \emph{vāda}, \textquotedblleft{}speech, dispute\textquotedblright{} (Monier-Williams), to \mr{घालणे}, \textquotedblleft{}to put on\textquotedblright{}.
\end{cousinweb}

\begin{grammarlens}[title={What you've built}]
One more word for giving reasons and opinions.
\end{grammarlens}

\subsection*{Your turn}
\begin{itemize}
\raggedright
  \item \emph{Say it:} \emph{vād ghālṇe}
  \item \emph{Say it:} \emph{vād ghālṇe}, once more
  \item \emph{Say it:} \emph{mulaṁ kheḷaṇyāṁvarūn vād ghālat hotī.} — The children were arguing over the toys.
  \item \emph{Recall:} say the Marathi for an activist, then the Marathi for to negotiate, then say \emph{vād ghālṇe} again
\end{itemize}

\begin{tcolorbox}[breakable,colback=teal!4,colframe=teal!35!black,title={Before you move on}]
What does \mr{वाद} \mr{घालणे} mean? (\textquotedblleft{}To argue\textquotedblright{}.) Say it once more.
\end{tcolorbox}
\section[\texorpdfstring{\mr{रेटणे} (reṭṇe)}{रेटणे (reṭṇe)}]{\mr{रेटणे} (reṭṇe) — to push through}

\label{lesson:MR-C365-retne}



\begin{quote}
Before the new one: say the Marathi for to negotiate, then the Marathi for to argue.
\end{quote}

\subsection*{You'll want to know: \mr{रेटणे}}
\textbf{\mr{रेटणे}} — \emph{reṭṇe} — \textquotedblleft{}to push through\textquotedblright{}.

\begin{quote}
\mr{विरोध} \mr{असूनही} \mr{त्याने} \mr{आपला} \mr{मुद्दा} \mr{रेटला}.
\end{quote}

\begin{quote}
\emph{virodh asūnhī tyāne āplā muddā reṭlā.}
\end{quote}

\begin{quote}
Despite opposition he pushed his point through.
\end{quote}

Past, with -\mr{ने} (he …): \mr{त्याने} \mr{मुद्दा} \mr{रेटला}, \emph{tyāne muddā reṭlā}.

Also to shove a load: \mr{गाडी} \mr{रेटणे}, to push a cart.

\begin{cousinweb}
\textbf{\mr{रेटणे}} (\emph{reṭṇe}) is a Marathi verb whose origin is not settled.
\end{cousinweb}

\begin{grammarlens}[title={What you've built}]
One more word for giving reasons and opinions.
\end{grammarlens}

\subsection*{Your turn}
\begin{itemize}
\raggedright
  \item \emph{Say it:} \emph{reṭṇe}
  \item \emph{Say it:} \emph{reṭṇe}, once more
  \item \emph{Say it:} \emph{virodh asūnhī tyāne āplā muddā reṭlā.} — Despite opposition he pushed his point through.
  \item \emph{Recall:} say the Marathi for to negotiate, then the Marathi for to argue, then say \emph{reṭṇe} again
\end{itemize}

\begin{tcolorbox}[breakable,colback=teal!4,colframe=teal!35!black,title={Before you move on}]
What does \mr{रेटणे} mean? (\textquotedblleft{}To push through\textquotedblright{}.) Say it once more.
\end{tcolorbox}
\section[\texorpdfstring{\mr{मिटवणे} (miṭavṇe)}{मिटवणे (miṭavṇe)}]{\mr{मिटवणे} (miṭavṇe) — to resolve (a dispute)}

\label{lesson:MR-C365-mitavne}



\begin{quote}
Before the new one: say the Marathi for to argue, then the Marathi for to push through.
\end{quote}

\subsection*{You'll want to know: \mr{मिटवणे}}
\textbf{\mr{मिटवणे}} — \emph{miṭavṇe} — \textquotedblleft{}to resolve (a dispute)\textquotedblright{}.

\begin{quote}
\mr{आजोबांनी} \mr{दोघांचं} \mr{भांडण} \mr{मिटवलं}.
\end{quote}

\begin{quote}
\emph{ājobānnī doghāñcaṁ bhāṇḍaṇ miṭavlaṁ.}
\end{quote}

\begin{quote}
Grandfather settled the quarrel between the two.
\end{quote}

Past, with -\mr{ने} (he …): \mr{त्याने} \mr{भांडण} \mr{मिटवलं}, \emph{tyāne bhāṇḍaṇ miṭavlaṁ}.

\mr{मिटणे}, to be settled: \mr{भांडण} \mr{मिटलं}.

\begin{cousinweb}
\textbf{\mr{मिटवणे}} (\emph{miṭavṇe}) is built on \mr{मिटणे}, \textquotedblleft{}to be wiped out\textquotedblright{}, as its causative; CDIAL lists Marathi \emph{miṭṇẽ}, \textquotedblleft{}to efface\textquotedblright{}, under Sanskrit \emph{mṛṣṭa}, \textquotedblleft{}rubbed, wiped\textquotedblright{} (CDIAL 10299).
\end{cousinweb}

\begin{grammarlens}[title={What you've built}]
One more word for giving reasons and opinions.
\end{grammarlens}

\subsection*{Your turn}
\begin{itemize}
\raggedright
  \item \emph{Say it:} \emph{miṭavṇe}
  \item \emph{Say it:} \emph{miṭavṇe}, once more
  \item \emph{Say it:} \emph{ājobānnī doghāñcaṁ bhāṇḍaṇ miṭavlaṁ.} — Grandfather settled the quarrel between the two.
  \item \emph{Recall:} say the Marathi for to argue, then the Marathi for to push through, then say \emph{miṭavṇe} again
\end{itemize}

\begin{tcolorbox}[breakable,colback=teal!4,colframe=teal!35!black,title={Before you move on}]
What does \mr{मिटवणे} mean? (\textquotedblleft{}To resolve (a dispute)\textquotedblright{}.) Say it once more.
\end{tcolorbox}
\section[\texorpdfstring{\mr{संघर्ष} \mr{करणे} (saṅgharṣa karṇe)}{संघर्ष करणे (saṅgharṣa karṇe)}]{\mr{संघर्ष} \mr{करणे} (saṅgharṣa karṇe) — to fight for (a cause); to struggle}

\label{lesson:MR-C365-sangharsa-karne}



\begin{quote}
Before the new one: say the Marathi for to push through, then the Marathi for to resolve.
\end{quote}

\subsection*{You'll want to know: \mr{संघर्ष} \mr{करणे}}
\textbf{\mr{संघर्ष} \mr{करणे}} — \emph{saṅgharṣa karṇe} — \textquotedblleft{}to fight for (a cause); to struggle\textquotedblright{}.

\begin{quote}
\mr{त्याने} \mr{शिक्षणासाठी} \mr{खूप} \mr{संघर्ष} \mr{केला}.
\end{quote}

\begin{quote}
\emph{tyāne śikṣaṇāsāṭhī khūp saṅgharṣa kelā.}
\end{quote}

\begin{quote}
He struggled hard for his education.
\end{quote}

Past, with -\mr{ने} (he …): \mr{त्याने} \mr{संघर्ष} \mr{केला}, \emph{tyāne saṅgharṣa kelā}.

X \mr{शी} \mr{संघर्ष}, a struggle with X.

\begin{cousinweb}
\textbf{\mr{संघर्ष} \mr{करणे}} (\emph{saṅgharṣa karṇe}) joins the tatsama \mr{संघर्ष}, Sanskrit \emph{saṁgharṣa}, \textquotedblleft{}rubbing together, rivalry\textquotedblright{} (Monier-Williams), to \mr{करणे}, \textquotedblleft{}to do\textquotedblright{}.
\end{cousinweb}

\begin{grammarlens}[title={What you've built}]
One more word for giving reasons and opinions.
\end{grammarlens}

\subsection*{Your turn}
\begin{itemize}
\raggedright
  \item \emph{Say it:} \emph{saṅgharṣa karṇe}
  \item \emph{Say it:} \emph{saṅgharṣa karṇe}, once more
  \item \emph{Say it:} \emph{tyāne śikṣaṇāsāṭhī khūp saṅgharṣa kelā.} — He struggled hard for his education.
  \item \emph{Recall:} say the Marathi for to push through, then the Marathi for to resolve, then say \emph{saṅgharṣa karṇe} again
\end{itemize}

\begin{tcolorbox}[breakable,colback=teal!4,colframe=teal!35!black,title={Before you move on}]
What does \mr{संघर्ष} \mr{करणे} mean? (\textquotedblleft{}To fight for (a cause); to struggle\textquotedblright{}.) Say it once more.
\end{tcolorbox}
\section[(dialogue)]{Looking back: from Making Laws to Settling Disputes}

\label{lesson:MR-R365-review-358-365}



\begin{quote}
Say the Marathi for a constitution and a policy.
\end{quote}

\subsection*{Your turn}
Say these aloud:

\begin{itemize}
\raggedright
  \item a constitution, a policy, to enforce, to cancel, to grant
  \item a lawyer, a lawsuit, a hearing, a verdict, a claim
  \item equality, discrimination, inequality, poverty, literacy
  \item revenue, a budget, a fund, a subsidy, economic
  \item a movement, a demonstration, a hunger strike, to protest, to boycott
  \item power, corruption, transparency, self-interest, to incite
  \item an institution, a committee, an assembly, a spokesperson, an activist
  \item to negotiate, to argue, to push through, to resolve, to fight for
\end{itemize}

\begin{tcolorbox}[breakable,colback=teal!4,colframe=teal!35!black,title={Before you move on}]
A constitution? (\textbf{\mr{संविधान}}, \emph{saṁvidhān}.) A policy? (\textbf{\mr{धोरण}}, \emph{dhoraṇ}.) To enforce? (\textbf{\mr{लागू} \mr{करणे}}, \emph{lāgū karṇe}.) To cancel? (\textbf{\mr{रद्द} \mr{करणे}}, \emph{radda karṇe}.) To grant? (\textbf{\mr{बहाल} \mr{करणे}}, \emph{bahāl karṇe}.) A lawyer? (\textbf{\mr{वकील}}, \emph{vakīl}.) A lawsuit? (\textbf{\mr{खटला}}, \emph{khaṭlā}.) A hearing? (\textbf{\mr{सुनावणी}}, \emph{sunāvṇī}.) A verdict? (\textbf{\mr{निवाडा}}, \emph{nivāḍā}.) A claim? (\textbf{\mr{दावा}}, \emph{dāvā}.) Equality? (\textbf{\mr{समानता}}, \emph{samāntā}.) Discrimination? (\textbf{\mr{भेदभाव}}, \emph{bhedbhāv}.) Inequality? (\textbf{\mr{विषमता}}, \emph{viṣamtā}.) Poverty? (\textbf{\mr{दारिद्र्य}}, \emph{dāridrya}.) Literacy? (\textbf{\mr{साक्षरता}}, \emph{sākṣartā}.) Revenue? (\textbf{\mr{महसूल}}, \emph{mahsūl}.) A budget? (\textbf{\mr{अर्थसंकल्प}}, \emph{arthasaṅkalpa}.) A fund? (\textbf{\mr{निधी}}, \emph{nidhī}.) A subsidy? (\textbf{\mr{अनुदान}}, \emph{anudān}.) Economic? (\textbf{\mr{आर्थिक}}, \emph{ārthik}.) A movement? (\textbf{\mr{आंदोलन}}, \emph{āndolan}.) A demonstration? (\textbf{\mr{निदर्शन}}, \emph{nidarśan}.) A hunger strike? (\textbf{\mr{उपोषण}}, \emph{upoṣaṇ}.) To protest? (\textbf{\mr{निषेध} \mr{करणे}}, \emph{niṣedh karṇe}.) To boycott? (\textbf{\mr{बहिष्कार} \mr{घालणे}}, \emph{bahiṣkār ghālṇe}.) Power? (\textbf{\mr{सत्ता}}, \emph{sattā}.) Corruption? (\textbf{\mr{भ्रष्टाचार}}, \emph{bhraṣṭācār}.) Transparency? (\textbf{\mr{पारदर्शकता}}, \emph{pārdarśaktā}.) Self-interest? (\textbf{\mr{स्वार्थ}}, \emph{svārtha}.) To incite? (\textbf{\mr{चिथावणे}}, \emph{cithāvṇe}.) An institution? (\textbf{\mr{संस्था}}, \emph{saṁsthā}.) A committee? (\textbf{\mr{समिती}}, \emph{samitī}.) An assembly? (\textbf{\mr{सभा}}, \emph{sabhā}.) A spokesperson? (\textbf{\mr{प्रवक्ता}}, \emph{pravaktā}.) An activist? (\textbf{\mr{कार्यकर्ता}}, \emph{kāryakartā}.) To negotiate? (\textbf{\mr{वाटाघाटी} \mr{करणे}}, \emph{vāṭāghāṭī karṇe}.) To argue? (\textbf{\mr{वाद} \mr{घालणे}}, \emph{vād ghālṇe}.) To push through? (\textbf{\mr{रेटणे}}, \emph{reṭṇe}.) To resolve? (\textbf{\mr{मिटवणे}}, \emph{miṭavṇe}.) To fight for? (\textbf{\mr{संघर्ष} \mr{करणे}}, \emph{saṅgharṣa karṇe}.)
\end{tcolorbox}
